\section{An operator with a uniform spectral gap}
\label{sec:spectrum}

We use the data of Proposition~\ref{prop:data} and the module constructed
above.  In particular, the induced intrinsic balls supplied by
Lemma~\ref{lem:overlap} have radius at most $\sigma(P)/16$, are trees,
and contain fewer than $\eta|P|$ vertices.  All copies below retain their
indices, and $|P|=n_i$ for a copy of type $i$.

\subsection{The reduced algebra and its evaluation representations}

The action of $\Gamma$ on the compact dual $\Omega$ of $V$ is
\[
 (g\omega)(v)=\omega(g^{-1}v).
\]
Write $\alpha_g(f)(\omega)=f(g^{-1}\omega)$ for the corresponding action
on $C(\Omega)$.  Fourier transformation on the discrete abelian group $V$
identifies $C_r^*(V)$ with $C(\Omega)$.  The regular representation of the
semidirect product, or equivalently regular induction followed by this
Fourier transformation, gives
\begin{equation}
 C_r^*(\Gsur)=C_r^*(V\rtimes\Gamma)
 \cong C(\Omega)\rtimes_{\alpha,r}\Gamma.
 \label{eq:semidirect-fourier}
\end{equation}
We use this identification throughout, and write $u_g$ for the canonical
unitary corresponding to $g\in\Gamma$.  Thus
$u_gfu_g^*=\alpha_g(f)$.

For $\omega\in\Omega$, the regular representation induced by evaluation
at $\omega$ acts on $\ell^2(\Gamma)$.  After inverting the usual group
index, its formulas are
\begin{equation}
 \begin{split}
 (\pi_\omega(f)\xi)(h)&=f(h^{-1}\omega)\xi(h),\qquad f\in C(\Omega),\\
 (\pi_\omega(u_g)\xi)(h)&=\xi(hg),\qquad g\in\Gamma.
 \end{split}
 \label{eq:regular-evaluations}
\end{equation}
The direct sum of all evaluations is a faithful representation of
$C(\Omega)$; inducing it regularly shows that
$\bigoplus_{\omega\in\Omega}\pi_\omega$ is faithful for the reduced
crossed product.  In particular, the norm of an element of
\eqref{eq:semidirect-fourier} is the supremum of its norms in this family.

Let $\tau_{\Gsur}$ denote the canonical trace, normalized by $\tau_{\Gsur}(1)=1$.
For every $b\in C_r^*(\Gsur)$,
\begin{equation}
 \tau_{\Gsur}(b)=\int_\Omega
 \langle\pi_\omega(b)\delta_1,\delta_1\rangle\,
 \mathrm{d}m_\Omega(\omega).
 \label{eq:trace-integral}
\end{equation}
Indeed, for a finite crossed-product sum $b=\sum_g F_g u_g$ the diagonal
coefficient at $1$ is $F_1(\omega)$, and its Haar integral is the
canonical group trace under Fourier transformation.  Such sums are dense.
The diagonal coefficients converge uniformly under norm approximation,
so the identity extends to every $b$.  This also proves that the
integrand in \eqref{eq:trace-integral} is continuous.

For $s\in S^{\pm1}$, let $f_s\in C(\Omega)$ be the clopen-set projection
\[
 f_s(\omega)=\mathbf 1_{\{s\in\mathcal J(\omega(v_*))\}}.
\]
Define
\begin{equation}
 T=dI-\sum_{s\in S^{\pm1}}f_su_sf_{s^{-1}}
 \quad\in C_r^*(\Gsur).
 \label{eq:operator}
\end{equation}
This element is self-adjoint, since taking adjoints interchanges the
terms indexed by $s$ and $s^{-1}$.  Each $f_s$ is a function of a single
finite-valued coordinate and hence is a finite Fourier sum of characters
of $\Omega$ indexed by elements of $V$.  Thus the construction takes
place in the ordinary reduced group $C^*$-algebra in
\eqref{eq:semidirect-fourier}.

\begin{definition}
For $\omega\in\Omega$, its \emph{enabled graph} $G_\omega$ has vertex
set $\Gamma$.  At $h\in\Gamma$ the allowed outgoing labels are
$\mathcal J(X_\omega(h))$.  The Cayley edge from $h$ to $hs$ is enabled
precisely when $s$ is allowed at $h$ and $s^{-1}$ is allowed at $hs$.
A copy $P$ is \emph{valid for $\omega$} if $X_\omega|_P=w_P$.
Let $U_\omega$ be the complement in $\Gamma$ of the union of all valid
copies.
\end{definition}

The graph $G_\omega$ is undirected and simple, and all its degrees are
at most $d$, because the decoder returns either a $d$-element set or the
empty set.  Since
\[
 (h^{-1}\omega)(v_*)=\omega(hv_*)=X_\omega(h),
\]
the formulas \eqref{eq:regular-evaluations} give
\begin{equation}
 \pi_\omega(T)=dI-\Adj(G_\omega).
 \label{eq:enabled-operator}
\end{equation}
The degree bound implies $\|\Adj(G_\omega)\|\le d$, so this operator
is bounded, positive, and at most $2dI$.

On a valid copy $P$, the allowed port sets are exactly $J_P(h)$.
All copy edges are therefore enabled, and their $d$ incidences at each
vertex exhaust its allowed labels.  As $P$ is connected, it is a whole
isolated component of $G_\omega$, with precisely its copy edges.
Moreover, distinct valid copies are disjoint.  To see this, two valid
copies meeting at a vertex would both be the entire same component.
If their indices were distinct, root $\widehat X$ at a shared group
vertex.  Their centers would both have distance one from the root, so
Lemma~\ref{lem:overlap} would place their entire common vertex set in
an intrinsic ball of the first copy containing fewer than
$\eta|P|<|P|$ vertices, a contradiction.  It follows also that no enabled
edge joins $U_\omega$ to its complement.

\subsection{Deficits on invalid copies and on arbitrary finite sets}

For $B\subseteq P$, let $e_{\omega,P}(B)$ be the number of enabled
copy edges with both endpoints in $B$.

\begin{lemma}[Deficit inside an invalid copy]
\label{lem:local-deficit}
If $P$ is invalid for $\omega$, then every subset $B\subseteq P$ satisfies
\begin{equation}
 d|B|-2e_{\omega,P}(B)\ge c|B|.
 \label{eq:local-deficit}
\end{equation}
\end{lemma}

\begin{proof}
Write $n=|P|$, and let
$M=\{h\in P:X_\omega(h)\ne w_P(h)\}$.
By Lemma~\ref{lem:marginals}, the restriction $X_\omega|_P$ belongs to
$C_P$, and $w_P\in C_P$ by Lemma~\ref{lem:codes}.  Their difference is
nonzero, since $P$ is invalid.  The primal-distance bound therefore gives
$|M|>\eta n$.

At a vertex of $M$, the decoder returns either the empty set or a
$d$-element set different from the prescribed $J_P(h)$.  In either case
at least one prescribed incidence is missing.  Consequently, writing
$b=|B|$, we have both
\[
 db-2e_{\omega,P}(B)\ge |M\cap B|
 \quad\text{and}\quad
 db-2e_{\omega,P}(B)\ge |\partial_P B|.
\]
If $n-b\ge\eta n/2$, then
$\min(b,n-b)\ge(\eta/2)b$.  The edge-expansion bound in
Proposition~\ref{prop:data} yields
\[
 |\partial_P B|\ge\kappa\min(b,n-b)
 \ge(\eta\kappa/2)b\ge cb.
\]
Otherwise,
\[
 |M\cap B|\ge |M|-(n-b)>\eta n/2
 \ge(\eta/2)b\ge cb.
\]
Here we used $c=(\eta/2)\min(\kappa,1)$.  The assertion for $B=\varnothing$
is included.
\end{proof}

\begin{lemma}[Deficit on the remaining graph]
\label{lem:global-deficit}
For every $\omega\in\Omega$ and every finite set $Y\subseteq U_\omega$,
let $A$ be the graph on $Y$ consisting of all its internal enabled edges.
Then
\begin{equation}
 d|Y|-2|E(A)|\ge(c/2)|Y|.
 \label{eq:global-deficit}
\end{equation}
\end{lemma}

\begin{proof}
The local deficits cannot simply be added over copies, since their
vertices and edges can overlap.  Instead, we bound the cycle rank
$b_1(A)=\dim Z_1(A;\F_2)$.  Euler's identity
\[
 |E(A)|=|Y|-\#\pi_0(A)+b_1(A)
\]
holds also for isolated vertices and the empty graph.  It shows that
the desired deficit follows from
\[
 b_1(A)\le\bigl(d/2-1-c/4\bigr)|Y|.
\]
We will obtain this bound by charging the ranks of successive copy
coordinates to pairwise disjoint subsets of $Y$.

\emph{Eliminating cycle coordinates.}
We regard the finite-dimensional space
$\mathcal L=Z_1(A;\F_2)$ as a subspace of $Z_1(X;\F_2)$.
Lemma~\ref{lem:cycle-sum} provides the algebraic direct-sum decomposition
\[
 Z_1(X;\F_2)=\bigoplus_P Z_1(P;\F_2),
\]
with coordinate maps $\operatorname{pr}_P$.  These are the coordinates
in the algebraic decomposition, not restrictions of edge chains to $P$:
a projected cycle can contain edges absent from the original cycle.
Only finitely many
coordinates can be nonzero on $\mathcal L$: choose a finite basis of
$\mathcal L$ and take the union of its finite coordinate supports.
Call these copies $P_1,\ldots,P_N$, ordered by nonincreasing distance
of their centers from a fixed group-vertex root $R$ in $\widehat X$.
Ties may be resolved arbitrarily.

For $1\le j\le N+1$, set
\[
 \mathcal L_j=
 \{z\in\mathcal L:\operatorname{pr}_{P_i}z=0\text{ for }i<j\}.
\]
Then $\mathcal L_1=\mathcal L$, $\mathcal L_{N+1}=0$, and rank-nullity
gives
\begin{equation}
 \dim\mathcal L=
 \sum_{j=1}^N\dim\operatorname{pr}_{P_j}(\mathcal L_j).
 \label{eq:cycle-elimination}
\end{equation}
If $\mathcal L=0$, this is understood as an empty sum.

\emph{Localizing projected cycles.}
Fix a step $j$, put $P=P_j$, and choose the induced intrinsic tree ball
$B_P$ given by Lemma~\ref{lem:overlap} for this root.  It contains every
vertex of $P$ shared with any of the other remaining copies
$P_{j+1},\ldots,P_N$.  Let
\[
 D_P=B_P\cup(A\cap P),
\]
where $A\cap P$ has vertex set $Y\cap P$ and the edges common to $A$
and $P$.  We claim that
\begin{equation}
 \operatorname{pr}_P(\mathcal L_j)\subseteq Z_1(D_P;\F_2).
 \label{eq:projected-cycle-support}
\end{equation}
Indeed, let an edge $e$ occur with nonzero coefficient in
$\operatorname{pr}_P z$ for $z\in\mathcal L_j$.  If $e$ does not belong
to $A$, its coefficient must cancel with a coefficient from another
remaining copy, because the total chain $z$ is supported in $A$.
Both endpoints of $e$ are consequently shared with that copy and lie
in $B_P$.  As the ball is induced in $P$, the edge $e$ itself lies in
$B_P$.  All other edges of $\operatorname{pr}_P z$ lie in $A\cap P$.
The projection is already a cycle in $P$, so the claimed inclusion
follows.  In particular,
\begin{equation}
 \dim\operatorname{pr}_P(\mathcal L_j)\le b_1(D_P).
 \label{eq:projected-cycle-rank}
\end{equation}

Figure~\ref{fig:cycle-charging} summarizes the bookkeeping in the next
step. The projected cycle may acquire edges in the overlap tree, but its
other edges remain enabled edges on $Y$. The tree can contribute to a
cycle only through multiple attachments of an exterior component.

\begin{figure}[ht]
\centering
\begin{tikzpicture}[font=\small,>=Stealth,vertex/.style={circle,fill=black,inner sep=1.5pt}]
\node[circle,draw,fill=blue!8,minimum size=15mm] (b) at (0,0) {$B_P$};
\node[ellipse,draw,fill=orange!10,minimum width=35mm,minimum height=18mm] (c1) at (4,0.8) {$C_1$};
\node[ellipse,draw,fill=orange!10,minimum width=30mm,minimum height=14mm] (c2) at (4,-1.3) {$C_2$};
\draw (b.25) to[bend left=12] node[above,sloped] {attachments} (c1.155);
\draw (b.0) to[bend right=10] (c1.195);
\draw (b.-35) -- (c2.160);
\node[align=center,text width=31mm] at (0,-1.65) {overlap tree\\contracted schematically};
\node[align=center,text width=43mm] at (4,-2.6) {exterior vertices\\charged only to $P$};
\node[align=center,text width=31mm] at (0,1.6) {all remaining copies\\meet $P$ inside $B_P$};
\end{tikzpicture}

\caption{The graph $D_P=B_P\cup(A\cap P)$ used to bound a projected
cycle coordinate. The tree $B_P$ is shown as a single region, and
$C_1,C_2$ represent components after deleting its vertices. Multiple
attachments contribute $l_C-1$ to the cycle rank, where $l_C$ is the
number of attachment edges of $C$. Large girth separates
their exterior endpoints. Only vertices outside $B_P$ are charged to
$P$, so the charges for successively eliminated copies are disjoint.
The drawing is schematic and specifies no distances.}
\label{fig:cycle-charging}
\end{figure}

\emph{Bounding exterior cycles and attachments.}
If $P$ is valid, $Y\cap P=\varnothing$, so $D_P=B_P$ is a tree and the
right side is zero.  Suppose from now on that $P$ is invalid, and write
$\sigma=\sigma(P)$ and $r$ for the radius of $B_P$.  Thus
$r\le\sigma/16$ and $\sigma>\max(16,20/c)$.

Consider the connected components $C$ of the graph obtained from $D_P$
by deleting the vertices of $B_P$.  Their vertex sets partition
$(Y\cap P)\setminus B_P$.  Let $l_C$ be the number of edges joining
$C$ to $B_P$.  Since $B_P$ is a connected tree, the Euler characteristic
formula for graph cycle rank gives
\begin{equation}
 b_1(D_P)=\sum_C\bigl(b_1(C)+\max(0,l_C-1)\bigr).
 \label{eq:tree-attachment-rank}
\end{equation}
This formula includes components with no attachment: such a component
contributes only its own cycle rank.

All edges of $C$ are enabled copy edges, and all enabled copy edges
between its vertices occur in $C$, by the definition of $A$.
If $C$ is cyclic, it has at least $\sigma$ vertices.  Applying
Lemma~\ref{lem:local-deficit} to its vertex set and using connectedness,
we obtain
\begin{equation}
 \begin{split}
 b_1(C)&=|E(C)|-|C|+1\\
 &\le\bigl(d/2-1-c/2+1/\sigma\bigr)|C|.
 \end{split}
 \label{eq:exterior-cycle-rank}
\end{equation}
The same bound holds for an acyclic $C$, because then $b_1(C)=0$ and
the displayed coefficient is positive.  Here $d\ge3$ and
$0<c\le\eta/2=1/200$.

We next bound the extra rank caused by attachments to the tree.
Choose two distinct attachment edges, and denote their endpoints in
$C$ by $v,v'$ and their endpoints in $B_P$ by $b,b'$.  A shortest path
in $C$ from $v$ to $v'$, these two edges, and the tree path from $b'$ to
$b$ form a simple cycle in $P$.  This remains true if one pair of
endpoints agrees, using the corresponding path of length zero; both
pairs cannot agree because $P$ is simple.  The tree path has length at
most $2r$, since both its endpoints lie in the intrinsic ball of
radius $r$.  Girth therefore gives
\begin{equation}
 \dist_C(v,v')\ge\sigma-2r-2>\sigma/2.
 \label{eq:attachment-separation}
\end{equation}
In particular, all the attachment endpoints in $C$ are distinct.

If $l_C\ge2$, the intrinsic balls in $C$ of radius
$\lfloor\sigma/4\rfloor$ centered at these endpoints are pairwise
disjoint, by \eqref{eq:attachment-separation}.  Each contains at least
$\lfloor\sigma/4\rfloor+1>\sigma/4$ vertices: take an initial segment
of a shortest path to any other attachment endpoint.  Thus
\begin{equation}
 \max(0,l_C-1)\le(4/\sigma)|C|.
 \label{eq:attachment-packing}
\end{equation}
For $l_C=0$ or $1$ the same inequality is immediate.

Combining \eqref{eq:tree-attachment-rank},
\eqref{eq:exterior-cycle-rank}, and \eqref{eq:attachment-packing}, and
using $5/\sigma<c/4$, yields
\begin{equation}
 \begin{split}
 b_1(D_P)
 &\le\bigl(d/2-1-c/2+5/\sigma\bigr)
       |(Y\cap P)\setminus B_P|\\
 &\le\bigl(d/2-1-c/4\bigr)|(Y\cap P)\setminus B_P|.
 \end{split}
 \label{eq:exterior-rank-bound}
\end{equation}
As observed above, this last bound holds also for a valid copy, since
both sides then vanish.

\emph{Summing over disjoint vertex sets.}
The exterior sets charged in \eqref{eq:exterior-rank-bound} are
pairwise disjoint.  Indeed, if $j<k$, the center of $P_k$ is no farther
from $R$ than the center of $P_j$.  Lemma~\ref{lem:overlap} therefore
implies
\[
 P_j\cap P_k\subseteq B_{P_j}.
\]
Consequently $(Y\cap P_j)\setminus B_{P_j}$ misses all of $P_k$,
including its charged exterior.  The non-strict distance comparison
also covers ties in the elimination order.

Summing \eqref{eq:projected-cycle-rank} and
\eqref{eq:exterior-rank-bound} in \eqref{eq:cycle-elimination} now gives
\[
 \dim Z_1(A;\F_2)\le\bigl(d/2-1-c/4\bigr)|Y|.
\]
The coefficient is positive, so this statement also follows from the
empty sum when $\mathcal L=0$.  Substituting this bound into Euler's
identity proves \eqref{eq:global-deficit}.
\end{proof}

\subsection{The spectral gap and its projection}

The passage from finite-set expansion to a spectral estimate belongs to
the discrete Cheeger tradition \cite{Dodziuk1984,AlonMilman1985}. We give
the coarea and Cauchy--Schwarz calculation to account for the vertex
potential $d-\deg_\omega(x)$ and the unnormalized Laplacian used here.

\begin{lemma}[Uniform spectral gap]
\label{lem:spectral-gap}
For every $\omega\in\Omega$, the restriction of $\pi_\omega(T)$ to
$\ell^2(U_\omega)$ satisfies
\begin{equation}
 \pi_\omega(T)|_{\ell^2(U_\omega)}\ge\frac{c^2}{8d}I.
 \label{eq:remaining-spectral-bound}
\end{equation}
Put
\[
 \varepsilon=\min\{\rho,c^2/(8d)\}>0.
\]
Then
\begin{equation}
 \spec(\pi_\omega(T))\subseteq\{0\}\cup[\varepsilon,2d],
 \label{eq:uniform-spectrum}
\end{equation}
and its kernel is exactly the Hilbert-space direct sum of the
one-dimensional spaces of constants on the valid copies.
\end{lemma}

\begin{proof}
On $U_\omega$ write $\deg_\omega(x)$ for the enabled degree and put
$q(x)=d-\deg_\omega(x)\ge0$.  There are no enabled edges to the removed
set.  Thus, for finitely supported complex $\xi$ on $U_\omega$,
\begin{equation}
 \begin{split}
 \mathcal E(\xi)
 :=\langle\pi_\omega(T)\xi,\xi\rangle
 =\sum_{\{x,y\}\in E(G_\omega|_{U_\omega})}
       |\xi(x)-\xi(y)|^2
   +\sum_{x\in U_\omega}q(x)|\xi(x)|^2.
 \end{split}
 \label{eq:remaining-energy}
\end{equation}
All edge sums in the rest of this proof are over enabled unoriented
edges on $U_\omega$.  For a finite set $F\subseteq U_\omega$, its
deficit is exactly
\[
 d|F|-2|E(G_\omega|_F)|
 =|\partial_{G_\omega|_{U_\omega}}F|+\sum_{x\in F}q(x).
\]

Set $s(x)=|\xi(x)|$.  Apply Lemma~\ref{lem:global-deficit} to the finite
level sets $\{x:s(x)^2>t\}$ and integrate in $t\ge0$.  The layer-cake
identity gives
\[
 (c/2)\|\xi\|^2
 \le\sum_{\{x,y\}}|s(x)^2-s(y)^2|+\sum_xq(x)s(x)^2.
\]
Cauchy--Schwarz, applied together to the edge terms and the
vertex-potential terms, then gives
\begin{align*}
 (c/2)\|\xi\|^2
 &\le
 \left(\sum_{\{x,y\}}|s(x)-s(y)|^2+\sum_xq(x)s(x)^2\right)^{1/2}
 \left(\sum_{\{x,y\}}(s(x)+s(y))^2+\sum_xq(x)s(x)^2\right)^{1/2}\\
 &\le\mathcal E(\xi)^{1/2}\sqrt{2d}\,\|\xi\|.
\end{align*}
For the first factor we used
$||\xi(x)|-|\xi(y)||\le|\xi(x)-\xi(y)|$.  For the second we used
\[
 \sum_{\{x,y\}}(s(x)+s(y))^2+\sum_xq(x)s(x)^2
 \le\sum_x\bigl(2\deg_\omega(x)+q(x)\bigr)s(x)^2
 \le2d\|\xi\|^2.
\]
If $\xi\ne0$, division and squaring prove
$\mathcal E(\xi)\ge(c^2/(8d))\|\xi\|^2$; the zero vector satisfies
the same inequality.  Finite-support vectors are dense and the operator
is bounded, so \eqref{eq:remaining-spectral-bound} holds on all of
$\ell^2(U_\omega)$.  If $U_\omega$ is empty, this assertion is vacuous.

On each valid copy $P$ of type $i$, the operator is its graph Laplacian
$dI-\Adj(P)$.  Its kernel consists of constants, and on the orthogonal
complement of constants it is at least $\rho I$, by
Proposition~\ref{prop:data}.  The valid copies are disjoint isolated
components.  Decomposing $\ell^2(\Gamma)$ into their component spaces
and $\ell^2(U_\omega)$ therefore shows that the full kernel is exactly
\[
 \bigoplus_{P\text{ valid for }\omega}\C\mathbf 1_P,
\]
and that the operator is at least $\varepsilon I$ on its orthogonal
complement.  These conclusions remain true for infinitely many valid
components, because both lower bounds are uniform.  The earlier bound
$\pi_\omega(T)\le2dI$ proves \eqref{eq:uniform-spectrum}.
\end{proof}

\begin{proposition}[A projection in the reduced group algebra]
\label{prop:projection}
There exists a projection $b_0\in C_r^*(\Gsur)$ such that, for every
$\omega\in\Omega$, $\pi_\omega(b_0)$ is the orthogonal projection
onto $\ker\pi_\omega(T)$.  Explicitly, it is projection onto constants
on each valid copy and is zero on $\ell^2(U_\omega)$.  In particular,
\begin{equation}
 \langle\pi_\omega(b_0)\delta_1,\delta_1\rangle
 =
 \begin{cases}
  1/n_i,&\text{if $1$ belongs to a valid copy of type $i$},\\
  0,&\text{if $1\in U_\omega$}.
 \end{cases}
 \label{eq:projection-diagonal}
\end{equation}
The valid copy in the first case is unique.
\end{proposition}

\begin{proof}
Take the faithful direct sum of the representations
$\pi_\omega$.  Lemma~\ref{lem:spectral-gap} gives the same positive
constant $\varepsilon$ in every summand, so its spectrum is contained
in the closed set $\{0\}\cup[\varepsilon,2d]$.  Faithfulness preserves
the spectrum of the self-adjoint element $T$; hence
\[
 \spec_{C_r^*(\Gsur)}(T)\subseteq\{0\}\cup[\varepsilon,2d].
\]
Choose a continuous real-valued function $\chi$ on $[0,2d]$ with
$\chi(0)=1$ and $\chi=0$ on $[\varepsilon,2d]$, and set
$b_0=\chi(T)$ by continuous functional calculus.  Since $\chi^2=\chi$
on the spectrum, $b_0$ is a projection in $C_r^*(\Gsur)$.  This uses the
ordinary C*-algebra functional calculus, with no passage to a von
Neumann algebra.

Functional calculus commutes with each $\pi_\omega$.  The kernel
description in Lemma~\ref{lem:spectral-gap} thus identifies
$\pi_\omega(b_0)$ as claimed, also when there are no valid copies or
infinitely many of them.  The normalized constant vector on a valid
copy of type $i$ is $n_i^{-1/2}\mathbf 1_P$, so the diagonal of its
rank-one projection at any vertex of $P$ is $1/n_i$.  At vertices of
$U_\omega$ the diagonal is zero.  Disjointness of the valid copies
proves uniqueness and \eqref{eq:projection-diagonal}.
\end{proof}
